% generado por el cuaderno; no editar
AA & 3 & \code{gc} (gulp call), \code{hm} (hoot call), \code{sc} (squeak call) \\
AC & 6 & \code{chc} (chitter call), \code{gc} (growl call), \code{cc} (contact call), \code{whc} (whinnie call), \code{sc} (squeak call), \code{bc} (bark call) \\
AS & 2 & \code{hc} (howl call), \code{bc} (bark call) \\
CC & 1 & \code{cc} (contact call) \\
LW & 11 & \code{cc} (contact call), \code{cs} (contact syllable), \code{aa} (aerial alarm call), \code{ta} (terrestrial alarm call), \code{sqc} (squeal call), \code{phc} (phee call), \code{tr} (trino call), \code{tf} (trino fast call), \code{tj} (unofficial tj call), \code{tt} (trino transition call), \code{vc} (visual contact call) \\
PT & 5 & \code{dc} (duet call), \code{ac} (alarm call), \code{bp} (unofficial bellow phrase), \code{pp} (unofficial pant phrase), \code{sqc} (unofficial squeal call) \\
SB & 5 & \code{pcc} (peep contact call), \code{ppc} (play peep call), \code{lpc} (long peep call), \code{spc} (spit call), \code{sc} (shriek call) \\
SM & 8 & \code{cc} (contact call), \code{acc} (aggressive contact call), \code{sc} (squeal call), \code{pc} (purr call), \code{whc} (whistle call), \code{hic} (hip call), \code{fc} (food call), \code{fs} (food syllable)
